\documentclass[11pt,a4paper]{scrartcl}

\usepackage[T1]{fontenc}
\usepackage[utf8]{inputenc}
\usepackage[ngerman]{babel}
\usepackage{lmodern}
\usepackage{microtype}
\usepackage{geometry}
\geometry{left=20mm,right=20mm,top=21mm,bottom=23mm}

\usepackage{booktabs}
\usepackage{longtable}
\usepackage{array}
\usepackage{ragged2e}
\usepackage{enumitem}
\usepackage{xcolor}
\usepackage{colortbl}
\usepackage[most]{tcolorbox}
\usepackage{pdflscape}
\usepackage{hyperref}
\usepackage{fancyhdr}
\usepackage{lastpage}
\usepackage{setspace}

\definecolor{TuringBlue}{HTML}{17365D}
\definecolor{TuringTeal}{HTML}{2F6F73}
\definecolor{TuringGreen}{HTML}{548235}
\definecolor{TuringOrange}{HTML}{C65911}
\definecolor{TuringSlate}{HTML}{5B6573}
\definecolor{TuringPurple}{HTML}{7030A0}

\definecolor{CategoryData}{HTML}{D9EAF7}
\definecolor{CategoryFunctional}{HTML}{E2F0D9}
\definecolor{CategoryArchitecture}{HTML}{E4DFEC}
\definecolor{CategoryAssurance}{HTML}{FCE4D6}
\definecolor{CategorySemantic}{HTML}{FFF2CC}

\newcommand{\statusdone}[1]{\textcolor{TuringGreen}{\textbf{#1}}}
\newcommand{\statuspartial}[1]{\textcolor{TuringOrange}{\textbf{#1}}}
\newcommand{\statusopen}[1]{\textcolor{TuringSlate}{\textbf{#1}}}
\newcommand{\statusfuture}[1]{\textcolor{TuringPurple}{\textbf{#1}}}

\hypersetup{
  colorlinks=true,
  linkcolor=TuringBlue,
  urlcolor=TuringTeal,
  pdftitle={Stakeholder Requirements und Implementierungsabdeckung},
  pdfauthor={Moritz Diez}
}

\setkomafont{section}{\color{TuringBlue}\Large\bfseries}
\setkomafont{subsection}{\color{TuringTeal}\large\bfseries}
\setlength{\parindent}{0pt}
\setlength{\parskip}{5pt plus 1pt minus 1pt}
\onehalfspacing

\pagestyle{fancy}
\fancyhf{}
\fancyhead[L]{\footnotesize Turing Generator -- Stakeholder Requirements}
\fancyhead[R]{\footnotesize Arbeitsstand 31.07.2026}
\fancyfoot[C]{\footnotesize Seite \thepage\ von \pageref{LastPage}}
\renewcommand{\headrulewidth}{0.4pt}

\begin{document}

\begin{titlepage}
  \centering
  \vspace*{2.0cm}
  {\color{TuringBlue}\rule{\textwidth}{1.2pt}}\\[1.2cm]
  {\Huge\bfseries\color{TuringBlue}
  Stakeholder Requirements und Implementierungsabdeckung\par}
  \vspace{0.7cm}
  {\Large Ausführliche Requirement-Tabelle des Turing Generators\par}
  \vspace{1.4cm}

  \begin{tcolorbox}[
    width=0.84\textwidth,
    colback=blue!3,
    colframe=TuringBlue,
    arc=2mm,
    boxrule=0.8pt
  ]
  \centering
  \textbf{Darstellungsperspektive}\par
  39 Stakeholder Requirements mit normativem Inhalt, Rationale, Quelle,\\
  Systemfunktionsabdeckung, Implementierungsstand und verbleibendem Zielumfang
  \end{tcolorbox}

  \vfill
  {\large Moritz Diez\par}
  \vspace{3mm}
  {\normalsize Masterarbeit Systems Engineering\par}
  \vspace{8mm}
  {\normalsize Arbeitsstand: 31. Juli 2026\par}
  \vspace{1.3cm}
  {\color{TuringBlue}\rule{\textwidth}{1.2pt}}
\end{titlepage}

\tableofcontents
\clearpage

% -----------------------------------------------------------------------------
% Inhaltsteil zum direkten Einbinden in die Masterarbeit
% Benoetigte Pakete:
% booktabs, longtable, array, ragged2e, xcolor, colortbl, tcolorbox,
% enumitem, pdflscape, hyperref, microtype
% -----------------------------------------------------------------------------

\section{Stakeholder Requirements und Implementierungsabdeckung}
\label{sec:stakeholder_requirements_coverage}

\subsection{Zweck der Requirement-Sicht}

Die Systemfunktionen beschreiben, welche fachlichen Leistungen der Turing
Generator erbringt. Die Stakeholder Requirements dokumentieren dagegen,
\emph{welche beobachtbaren Ergebnisse und Qualitäten aus Sicht der Stakeholder
erwartet werden}. Die nachfolgende Requirement-Tabelle stellt deshalb die
normative Stakeholder-Sicht in den Mittelpunkt.

Insgesamt umfasst die aktuelle CATIA-Baseline 39 Stakeholder Requirements. Sie
sind über die System Requirements vollständig mit den zwölf Systemfunktionen
verknüpft. Die detaillierten System Requirements werden in dieser Übersicht
bewusst nicht eingeblendet. Die sichtbare Traceability lautet damit:

\begin{center}
\textbf{Stakeholder Requirement}
$\longrightarrow$
\textbf{System Requirement}
$\longrightarrow$
\textbf{System Function}
$\longrightarrow$
\textbf{Implementierungsevidenz}
\end{center}

\begin{tcolorbox}[
  enhanced,
  breakable,
  colback=blue!3,
  colframe=TuringBlue,
  title={Lesart der Tabelle},
  fonttitle=\bfseries,
  boxrule=0.8pt,
  arc=2mm
]
Die Spalte \emph{Systemfunktionen} zeigt die fachliche Modellabdeckung eines
Stakeholder Requirements. Die Spalte \emph{Aktueller Stand} bewertet dagegen
die reale Implementierung in der Phase-F/P-Baseline. Vollständige
Modellabdeckung ist daher nicht gleichbedeutend mit vollständiger
Implementierung.
\end{tcolorbox}

\subsection{Kategorien und Statusdefinitionen}

Die 39 Stakeholder Requirements sind in fünf Anforderungskategorien gegliedert.

\begin{table}[htbp]
\centering
\caption{Kategorien der Stakeholder Requirements}
\label{tab:sr_categories}
\begin{tabular}{p{3.3cm}p{1.7cm}p{9.4cm}}
\toprule
\textbf{Kategorie} & \textbf{Anzahl} & \textbf{Schwerpunkt} \\
\midrule
Datenqualität & 8 & Eignung, Vollständigkeit, Konsistenz, Constraints, Completion und Evidenzabdeckung der Eingangsinformationen. \\
Funktional & 9 & Fachliche Systemleistungen von Projekt- und Quellenverwaltung bis zur SysML-v2-Zielrepräsentation. \\
Architektur & 10 & Strukturkonformität, Modularität, Wiederverwendung, Integration, Constraints, Patterns und Schnittstellen. \\
Absicherung & 9 & Ambiguitäten, Erklärbarkeit, Human Review, Validierung, Freigabegrenzen und Verarbeitungshistorie. \\
Semantik & 3 & Technisches Vokabular, semantische Äquivalenz und Ontologiekonflikte. \\
\bottomrule
\end{tabular}
\end{table}

\begin{table}[htbp]
\centering
\caption{Bedeutung der verwendeten Implementierungsstatus}
\label{tab:sr_status_definitions}
\begin{tabular}{p{4.1cm}p{10.3cm}}
\toprule
\textbf{Status} & \textbf{Bedeutung} \\
\midrule
\statusdone{Implementiert und verifiziert} & Das Stakeholder Requirement ist innerhalb der aktuellen Systemgrenze durch implementierte Funktionalität und Test- oder Review-Evidenz belastbar abgedeckt. \\
\statuspartial{Teilweise implementiert} & Wesentliche Grundlagen oder Teilfähigkeiten existieren, der vollständige normative Zielumfang ist jedoch noch nicht erreicht. \\
\statusopen{Architektur definiert} & Requirement, Systemfunktion und Architekturverantwortung sind modelliert; die zentrale Runtime-Fähigkeit ist noch nicht implementiert. \\
\statusfuture{Optionaler Ausblick} & Das Requirement liegt ausdrücklich außerhalb des aktuellen MVP und benötigt vor einer Umsetzung eine eigene Architektur- und Validierungsentscheidung. \\
\bottomrule
\end{tabular}
\end{table}

\subsection{Referenz der Systemfunktionen}

Die Requirement-Tabelle verwendet aus Platzgründen die IDs der zwölf
Systemfunktionen. Die folgende Legende dient als Kurzreferenz.

\begin{table}[htbp]
\centering
\caption{Systemfunktionen als Referenz für die Requirement-Tabelle}
\label{tab:sr_function_reference}
\begin{tabular}{p{1.7cm}p{6.0cm}p{1.7cm}p{6.0cm}}
\toprule
\textbf{ID} & \textbf{Kurzbeschreibung} & \textbf{ID} & \textbf{Kurzbeschreibung} \\
\midrule
SF\_001 & Projekt- und Quellenkontext verwalten & SF\_007 & Human Review und Freigabe unterstützen \\
SF\_002 & Quellinformationen bewerten und vorbereiten & SF\_008 & Architekturinformation ableiten und strukturieren \\
SF\_003 & Semantische Leitplanken anwenden & SF\_009 & Architekturinhalt validieren und integrieren \\
SF\_004 & Coverage und Evidenz bewerten & SF\_010 & SysML-v2-Artefakte erzeugen und validieren \\
SF\_005 & Processing Runs konfigurieren und steuern & SF\_011 & Evidenz und Traceability erhalten \\
SF\_006 & Engineering Candidates erzeugen und vergleichen & SF\_012 & Status und Historie darstellen \\
\bottomrule
\end{tabular}
\end{table}

\clearpage
\pagestyle{empty}
\begin{landscape}
\subsection{Ausführliche Stakeholder-Requirement-Tabelle}
\label{sec:detailed_sr_table}

\setlength{\LTleft}{0pt}
\setlength{\LTright}{0pt}
\setlength{\tabcolsep}{2.1pt}
\renewcommand{\arraystretch}{1.12}
\scriptsize

\begin{longtable}{
  >{\RaggedRight\arraybackslash}p{1.25cm}
  >{\RaggedRight\arraybackslash}p{4.65cm}
  >{\RaggedRight\arraybackslash}p{1.90cm}
  >{\RaggedRight\arraybackslash}p{4.00cm}
  >{\RaggedRight\arraybackslash}p{2.55cm}
  >{\RaggedRight\arraybackslash}p{5.05cm}
  >{\RaggedRight\arraybackslash}p{4.05cm}
}
\caption{Ausführliche Stakeholder-Requirement-Tabelle mit Funktions- und Implementierungsabdeckung}
\label{tab:detailed_stakeholder_requirements}\\
\toprule
\textbf{ID} &
\textbf{Stakeholder Requirement} &
\textbf{Kategorie} &
\textbf{Rationale und Quelle} &
\textbf{\shortstack[l]{System-\\funktionen}} &
\textbf{Aktueller Stand} &
\textbf{Verbleibender Zielumfang} \\
\midrule
\endfirsthead

\multicolumn{7}{l}{\footnotesize\itshape Fortsetzung von Tabelle~\ref{tab:detailed_stakeholder_requirements}}\\
\toprule
\textbf{ID} &
\textbf{Stakeholder Requirement} &
\textbf{Kategorie} &
\textbf{Rationale und Quelle} &
\textbf{\shortstack[l]{System-\\funktionen}} &
\textbf{Aktueller Stand} &
\textbf{Verbleibender Zielumfang} \\
\midrule
\endhead

\midrule
\multicolumn{7}{r}{\footnotesize Fortsetzung auf der nächsten Seite}\\
\endfoot

\bottomrule
\endlastfoot

\multicolumn{7}{l}{\cellcolor{CategoryData}\textbf{Datenqualität} -- Bewertung der Eignung, Vollständigkeit, Konsistenz und Evidenz der Eingangsinformationen.} \\
\addlinespace[1.5pt]
\textbf{SR\_001} & \textbf{Input Data Suitability Assessment}\par\smallskip The system shall assess whether legacy engineering data is suitable for automated processing. & Datenqualität & \textbf{Rationale:} Im Stakeholder Requirement nicht separat dokumentiert.\par\smallskip \textbf{Source:} \texttt{UN\_01\_001; UN\_01\_002; UN\_04\_001} & SF\_002, SF\_005 & \statusdone{Implementiert und verifiziert}\par\smallskip Phase F bewertet die Verarbeitbarkeit registrierter Quellen; P9 bindet die Prüfung an Projekt, Source Projection und Processing Run. & Eignungskriterien für zusätzliche Quelltypen und spätere modellbildende Stufen weiter operationalisieren. \\
\addlinespace[2pt]
\textbf{SR\_002} & \textbf{Input Data Completeness Check}\par\smallskip The system shall identify incomplete input data before it is used for downstream engineering processing. & Datenqualität & \textbf{Rationale:} Im Stakeholder Requirement nicht separat dokumentiert.\par\smallskip \textbf{Source:} \texttt{UN\_01\_004; UN\_04\_003} & SF\_002, SF\_005 & \statusdone{Implementiert und verifiziert}\par\smallskip Gaps und fehlende Informationen werden vor der nachgelagerten Nutzung im Ingestion Review Report ausgewiesen. & Vollständigkeitskriterien für spätere Architektur- und Generierungsstufen ergänzen. \\
\addlinespace[2pt]
\textbf{SR\_003} & \textbf{Input Data Consistency Check}\par\smallskip The system shall identify inconsistent input data before it is used for downstream engineering processing. & Datenqualität & \textbf{Rationale:} Im Stakeholder Requirement nicht separat dokumentiert.\par\smallskip \textbf{Source:} \texttt{UN\_01\_030; UN\_04\_003} & SF\_002, SF\_005 & \statuspartial{Teilweise implementiert}\par\smallskip Inkonsistenzen erscheinen als Risks, Ambiguities und Review Findings; eine vollständige quellenübergreifende Konsistenzprüfung ist noch nicht geschlossen. & Deterministische Konsistenzregeln über mehrere Quellen und spätere Modellartefakte implementieren. \\
\addlinespace[2pt]
\textbf{SR\_004} & \textbf{Source Reliability Classification}\par\smallskip The system shall classify source information according to defined reliability criteria. & Datenqualität & \textbf{Rationale:} Im Stakeholder Requirement nicht separat dokumentiert.\par\smallskip \textbf{Source:} \texttt{UN\_01\_003; UN\_04\_002} & SF\_002, SF\_005 & \statuspartial{Teilweise implementiert}\par\smallskip Quellenmetadaten, Rollen und Review-Evidence existieren; eine eigenständige, vollständig operationalisierte Reliability-Klassifikation ist nur teilweise abgebildet. & Reliability-Kriterien, Skala und nachvollziehbare Klassifikationsentscheidung vervollständigen. \\
\addlinespace[2pt]
\textbf{SR\_005} & \textbf{Source Relevance Classification}\par\smallskip The system shall classify source information according to its relevance to the selected engineering task. & Datenqualität & \textbf{Rationale:} Im Stakeholder Requirement nicht separat dokumentiert.\par\smallskip \textbf{Source:} \texttt{UN\_01\_001; UN\_04\_012} & SF\_001, SF\_002, SF\_005 & \statuspartial{Teilweise implementiert}\par\smallskip Engineering-source/context-only-Rollen, Framework-Zuordnung und taskbezogene Kontextselektion bilden wesentliche Relevanzmechanismen. & Explizite Relevanzklassifikation für jede ausgewählte Engineering Task konsolidieren. \\
\addlinespace[2pt]
\textbf{SR\_023} & \textbf{Input Constraint Violation Detection}\par\smallskip The system shall identify violations of defined input constraints in legacy engineering data. & Datenqualität & \textbf{Rationale:} Im Stakeholder Requirement nicht separat dokumentiert.\par\smallskip \textbf{Source:} \texttt{UN\_01\_005} & SF\_002 & \statuspartial{Teilweise implementiert}\par\smallskip Phase F erkennt Risks und Verletzungen; einzelne technische Verträge und Manifestregeln werden fail-closed validiert. & Einheitlichen Input-Constraint-Katalog für Engineering-Inhalte definieren. \\
\addlinespace[2pt]
\textbf{SR\_026} & \textbf{Legacy Data Completion}\par\smallskip The system shall support controlled completion of incomplete legacy engineering data before downstream engineering processing. & Datenqualität & \textbf{Rationale:} Incomplete legacy information may require reviewed domain knowledge before it can be used for subsequent engineering processing.\par\smallskip \textbf{Source:} \texttt{UN\_01\_004; UN\_02\_003; UN\_02\_006} & SF\_002, SF\_011 & \statuspartial{Teilweise implementiert}\par\smallskip Review Decisions, Terminology Decisions und persistente Evidenz bilden Grundlagen für kontrollierte Ergänzungen. & Expliziten Completion Workflow mit Domain Knowledge, Review und Approved Input implementieren. \\
\addlinespace[2pt]
\textbf{SR\_033} & \textbf{Engineering Information Coverage}\par\smallskip The system shall indicate which defined model areas are supported by available engineering information. & Datenqualität & \textbf{Rationale:} Indicated information coverage shall not be interpreted as approved model-generation readiness.\par\smallskip \textbf{Source:} \texttt{UN\_01\_027} & SF\_004, SF\_011 & \statusdone{Implementiert und verifiziert}\par\smallskip Preliminary Coverage unterscheidet covered, uncovered und attention-required Framework Nodes und stellt die Evidenz dar. & Approved Generation Readiness getrennt ergänzen und evaluieren. \\
\addlinespace[2pt]
\multicolumn{7}{l}{\cellcolor{CategoryFunctional}\textbf{Funktional} -- Fachliche Leistungen des Systems von Projektverwaltung bis Architecture as Code.} \\
\addlinespace[1.5pt]
\textbf{SR\_006} & \textbf{Architecture Information Derivation}\par\smallskip The system shall derive system-architecture content from heterogeneous legacy engineering information. & Funktional & \textbf{Rationale:} Separates the derivation of system-architecture content from its textual representation and the required target notation.\par\smallskip \textbf{Source:} \texttt{UN\_01\_006} & SF\_008 & \statusopen{Architektur definiert}\par\smallskip RFLP-Zielstruktur, Mapping Targets und Logical Architecture sind definiert; die eigentliche Architekturableitung ist noch nicht implementiert. & Architecture Candidate Layer und automatisierte Derivation umsetzen. \\
\addlinespace[2pt]
\textbf{SR\_019} & \textbf{Source Traceability}\par\smallskip The system shall maintain traceability between each generated engineering result and its originating source information. & Funktional & \textbf{Rationale:} Im Stakeholder Requirement nicht separat dokumentiert.\par\smallskip \textbf{Source:} \texttt{UN\_01\_019; UN\_01\_020; UN\_04\_006} & SF\_001, SF\_002, SF\_010, SF\_011 & \statusdone{Implementiert und verifiziert}\par\smallskip Source IDs, Locators, Hashes, Information Units, Processing Runs und Artifact References erhalten die Herkunft durchgängig. & Traceability bis zu finalen SysML-v2-Elementen und CATIA-Objekten verlängern. \\
\addlinespace[2pt]
\textbf{SR\_027} & \textbf{Architecture as Code Provision}\par\smallskip The system shall provide system architecture as machine-processable textual artifacts. & Funktional & \textbf{Rationale:} Enables system architecture to be stored, processed, compared and transferred as textual engineering information.\par\smallskip \textbf{Source:} \texttt{UN\_01\_022} & SF\_010 & \statusopen{Architektur definiert}\par\smallskip Architecture-as-Code ist als Zielrepräsentation und versionierte Artefaktstruktur definiert. & Maschinenverarbeitbare Architekturartefakte zur Laufzeit erzeugen. \\
\addlinespace[2pt]
\textbf{SR\_028} & \textbf{SysML v2 Target Representation}\par\smallskip The system shall represent architecture artifacts in SysML v2 textual notation. & Funktional & \textbf{Rationale:} SysML v2 is the target architecture representation defined by the scope of the Master's thesis.\par\smallskip \textbf{Source:} \texttt{UN\_01\_023; Master's thesis proposal} & SF\_010 & \statusopen{Architektur definiert}\par\smallskip SysML v2 Target Notation und CATIA/SYSIDE-Zielumgebung sind festgelegt. & Automatische SysML-v2-Erzeugung und Zielumgebungsvalidierung implementieren. \\
\addlinespace[2pt]
\textbf{SR\_029} & \textbf{Multiple Project Support}\par\smallskip The system shall maintain multiple distinct engineering projects. & Funktional & \textbf{Rationale:} Multiple projects may coexist within the system. This requirement does not require parallel processing execution.\par\smallskip \textbf{Source:} \texttt{UN\_01\_024} & SF\_001 & \statusdone{Implementiert und verifiziert}\par\smallskip Mehrere persistente Projekt-Workspaces können angelegt, ausgewählt, geladen und deterministisch wieder geöffnet werden. & Projektbearbeitung und kontrolliertes Löschen ergänzen. \\
\addlinespace[2pt]
\textbf{SR\_031} & \textbf{Processing State Transparency}\par\smallskip The system shall indicate the current processing state of each registered engineering source. & Funktional & \textbf{Rationale:} Im Stakeholder Requirement nicht separat dokumentiert.\par\smallskip \textbf{Source:} \texttt{UN\_01\_025} & SF\_005, SF\_011, SF\_012 & \statusdone{Implementiert und verifiziert}\par\smallskip Run- und Source-Aggregation, Event Chains und Dashboard zeigen den aktuellen Verarbeitungszustand. & Spätere Architektur- und Generierungszustände integrieren. \\
\addlinespace[2pt]
\textbf{SR\_032} & \textbf{Review Need Transparency}\par\smallskip The system shall identify processing results awaiting human review. & Funktional & \textbf{Rationale:} Im Stakeholder Requirement nicht separat dokumentiert.\par\smallskip \textbf{Source:} \texttt{UN\_01\_026} & SF\_011, SF\_012 & \statusdone{Implementiert und verifiziert}\par\smallskip review\_requested-Events, awaiting\_review und Attention/Review-Ansichten kennzeichnen offenen Review-Bedarf. & Mehrstufige Review-Bedarfe für Approved Input und Model Candidates ergänzen. \\
\addlinespace[2pt]
\textbf{SR\_037} & \textbf{Non-Textual Engineering Information}\par\smallskip The system shall extract engineering information from supported non-textual source artifacts. & Funktional & \textbf{Rationale:} Optional future capability outside the current MVP. Technical drawings are an example of a potential source type. The extraction and validation workflow is not yet defined.\par\smallskip \textbf{Source:} \texttt{UN\_04\_010; post-MVP outlook} & SF\_002, SF\_011 & \statusfuture{Optionaler Ausblick}\par\smallskip Die aktuelle MVP-Grenze unterstützt nur native Textquellen und deterministische Textprojektionen; PDF ist auf extrahierbare Textschichten begrenzt. & Separate Architektur für OCR, Bildverständnis und technische Zeichnungen definieren und validieren. \\
\addlinespace[2pt]
\textbf{SR\_038} & \textbf{Multiple Source Support}\par\smallskip The system shall maintain multiple distinct engineering source artifacts within each engineering project. & Funktional & \textbf{Rationale:} Im Stakeholder Requirement nicht separat dokumentiert.\par\smallskip \textbf{Source:} \texttt{UN\_01\_029} & SF\_001, SF\_002, SF\_005 & \statusdone{Implementiert und verifiziert}\par\smallskip Jedes Projekt kann mehrere getrennte Source Manifests und Source Artifacts verwalten. & Gemeinsame Multi-source-Runs und quellenübergreifende Verarbeitung ergänzen. \\
\addlinespace[2pt]
\multicolumn{7}{l}{\cellcolor{CategoryArchitecture}\textbf{Architektur} -- Strukturelle Eigenschaften, Modularität, Integrationsfähigkeit und Konformität der Zielarchitektur.} \\
\addlinespace[1.5pt]
\textbf{SR\_007} & \textbf{Structural Pattern Conformance}\par\smallskip The system shall generate model content according to the defined target-model structure. & Architektur & \textbf{Rationale:} This requirement addresses conformance to the mandatory target-model structure rather than the application of reusable architectural patterns.\par\smallskip \textbf{Source:} \texttt{UN\_01\_007} & SF\_008, SF\_009, SF\_010 & \statusopen{Architektur definiert}\par\smallskip Versionierte Framework- und Target-Structure-Artefakte definieren die Zielstruktur; es wird noch kein Modellinhalt generiert. & Konformitätsprüfung bei der späteren Architekturgenerierung implementieren. \\
\addlinespace[2pt]
\textbf{SR\_008} & \textbf{Model Reusability}\par\smallskip The system shall generate model elements that are reusable across modeling contexts. & Architektur & \textbf{Rationale:} Im Stakeholder Requirement nicht separat dokumentiert.\par\smallskip \textbf{Source:} \texttt{UN\_01\_008} & SF\_008 & \statusopen{Architektur definiert}\par\smallskip Wiederverwendbarkeit ist in Anforderungen und Zielarchitektur verankert, aber noch nicht an generierten Model Units nachgewiesen. & Wiederverwendbare Model Units erzeugen und in mehreren Modellierungskontexten evaluieren. \\
\addlinespace[2pt]
\textbf{SR\_009} & \textbf{Modular Model Growth}\par\smallskip The system shall preserve modular model structures during iterative model growth. & Architektur & \textbf{Rationale:} Im Stakeholder Requirement nicht separat dokumentiert.\par\smallskip \textbf{Source:} \texttt{UN\_01\_010; UN\_03\_001} & SF\_008 & \statusopen{Architektur definiert}\par\smallskip Modulare Zielstruktur und immutable Artefaktprinzipien sind definiert; iteratives Modellwachstum wird noch nicht ausgeführt. & Versioniertes Wachstum erzeugter Modelle und Regressionstests umsetzen. \\
\addlinespace[2pt]
\textbf{SR\_010} & \textbf{Non Destructive Submodel Integration}\par\smallskip The system shall enable generated submodels to be integrated without restructuring existing model elements. & Architektur & \textbf{Rationale:} Im Stakeholder Requirement nicht separat dokumentiert.\par\smallskip \textbf{Source:} \texttt{UN\_01\_009; UN\_03\_002} & SF\_008, SF\_009 & \statusopen{Architektur definiert}\par\smallskip Non-destructive Integration ist als Zielanforderung und Validierungsfunktion modelliert. & Integration erzeugter Submodelle in bestehende Kontexte implementieren und evaluieren. \\
\addlinespace[2pt]
\textbf{SR\_020} & \textbf{Integrated Model Consistency}\par\smallskip The system shall identify architectural inconsistencies between integrated model elements. & Architektur & \textbf{Rationale:} Im Stakeholder Requirement nicht separat dokumentiert.\par\smallskip \textbf{Source:} \texttt{UN\_01\_021; UN\_03\_004} & SF\_009 & \statusopen{Architektur definiert}\par\smallskip Die Validierungsfunktion für integrierte Architektur ist im Systemmodell definiert. & Modellübergreifende Konsistenzregeln und Integrationsprüfung implementieren. \\
\addlinespace[2pt]
\textbf{SR\_021} & \textbf{Larger Context Compatibility}\par\smallskip The system shall ensure synthesized model structures remain compatible with larger system contexts. & Architektur & \textbf{Rationale:} Im Stakeholder Requirement nicht separat dokumentiert.\par\smallskip \textbf{Source:} \texttt{UN\_03\_005} & SF\_009 & \statusopen{Architektur definiert}\par\smallskip Larger-context compatibility ist als explizite Zielanforderung enthalten. & Kontextschnittstelle, Compatibility Checks und Evaluationsszenarien implementieren. \\
\addlinespace[2pt]
\textbf{SR\_022} & \textbf{Architectural Constraint Conformance}\par\smallskip The system shall reflect defined architectural constraints in generated model content. & Architektur & \textbf{Rationale:} Im Stakeholder Requirement nicht separat dokumentiert.\par\smallskip \textbf{Source:} \texttt{UN\_03\_006} & SF\_008, SF\_009 & \statusopen{Architektur definiert}\par\smallskip System Design Constraints und Architekturvalidierungsfunktion sind modelliert. & Constraints maschinenlesbar auf generierten Architekturinhalt anwenden. \\
\addlinespace[2pt]
\textbf{SR\_024} & \textbf{Interface Definition Consistency}\par\smallskip The system shall maintain consistent interface definitions across generated model structures. & Architektur & \textbf{Rationale:} Im Stakeholder Requirement nicht separat dokumentiert.\par\smallskip \textbf{Source:} \texttt{UN\_03\_003} & SF\_008, SF\_009 & \statusopen{Architektur definiert}\par\smallskip Interface Consistency ist normativ verankert und der Logical Architecture zugeordnet. & Schnittstellen ableiten, vergleichen und konfliktfrei integrieren. \\
\addlinespace[2pt]
\textbf{SR\_025} & \textbf{Reusable Architectural Pattern Application}\par\smallskip The system shall apply selected reusable architectural patterns to generated model content. & Architektur & \textbf{Rationale:} This requirement addresses reusable architectural patterns and is distinct from mandatory target-model structure conformance.\par\smallskip \textbf{Source:} \texttt{UN\_03\_007} & SF\_008, SF\_009 & \statusopen{Architektur definiert}\par\smallskip Framework Templates und Mapping Targets schaffen die Grundlage für wiederverwendbare Muster. & Ausgewählte Patterns auf Model Candidates anwenden und ihre Konformität validieren. \\
\addlinespace[2pt]
\textbf{SR\_030} & \textbf{Project Information Isolation}\par\smallskip The system shall prevent engineering information assigned to one project from being used within another project. & Architektur & \textbf{Rationale:} Im Stakeholder Requirement nicht separat dokumentiert.\par\smallskip \textbf{Source:} \texttt{UN\_04\_007} & SF\_001, SF\_005, SF\_011 & \statusdone{Implementiert und verifiziert}\par\smallskip Projektlokale IDs, Pfade, Repositories und Cross-project-Rejection verhindern Informationsvermischung. & Isolationsnachweis auf spätere Approved-Input- und Modellartefakte erweitern. \\
\addlinespace[2pt]
\multicolumn{7}{l}{\cellcolor{CategoryAssurance}\textbf{Absicherung} -- Erklärbarkeit, Review, Validierung, Freigabegrenzen und nachvollziehbare Entscheidungen.} \\
\addlinespace[1.5pt]
\textbf{SR\_011} & \textbf{Ambiguity Detection}\par\smallskip The system shall identify ambiguities during automated processing of legacy engineering information. & Absicherung & \textbf{Rationale:} Im Stakeholder Requirement nicht separat dokumentiert.\par\smallskip \textbf{Source:} \texttt{UN\_01\_011; UN\_02\_001} & SF\_003, SF\_006 & \statusdone{Implementiert und verifiziert}\par\smallskip Phase F identifiziert Ambiguities; Multi-Agenten-Konsens erhält Disagreement und Varianz als Review-Evidence. & Erkennung auf spätere Architektur- und Modellkandidaten ausweiten. \\
\addlinespace[2pt]
\textbf{SR\_012} & \textbf{Expert Review of Generated Interpretations}\par\smallskip The system shall provide generated interpretations for expert review before their subsequent engineering use. & Absicherung & \textbf{Rationale:} Im Stakeholder Requirement nicht separat dokumentiert.\par\smallskip \textbf{Source:} \texttt{UN\_01\_012; UN\_02\_002} & SF\_005, SF\_007, SF\_012 & \statuspartial{Teilweise implementiert}\par\smallskip Review Reports, review\_requested-Events und Dashboard-Navigation stellen Interpretationen zur Prüfung bereit. & Durchgängigen Review bis Approved Input und später zur Modellfreigabe schließen. \\
\addlinespace[2pt]
\textbf{SR\_013} & \textbf{Human Intervention for Conflicts}\par\smallskip The system shall allow users to intervene when conflicting interpretations of source data occur. & Absicherung & \textbf{Rationale:} Im Stakeholder Requirement nicht separat dokumentiert.\par\smallskip \textbf{Source:} \texttt{UN\_01\_013; UN\_02\_003} & SF\_006, SF\_007 & \statuspartial{Teilweise implementiert}\par\smallskip Human Review Decisions können Ergebnisse auswählen, ablehnen oder dispositionieren und sind fingerprintgebunden. & Bedienbaren Konfliktlösungs- und Änderungsworkflow über alle Prozessstufen ergänzen. \\
\addlinespace[2pt]
\textbf{SR\_014} & \textbf{Explanation of Generated Engineering Results}\par\smallskip The system shall provide the rationale for generated engineering results. & Absicherung & \textbf{Rationale:} Im Stakeholder Requirement nicht separat dokumentiert.\par\smallskip \textbf{Source:} \texttt{UN\_01\_014} & SF\_002, SF\_003, SF\_004, SF\_006, SF\_010, SF\_012 & \statuspartial{Teilweise implementiert}\par\smallskip Agent Outputs, Konsensberichte, Review Reports und Evidenzreferenzen liefern Begründungen für aktuelle Verarbeitungsergebnisse. & Erklärbarkeit auf Architekturentscheidungen und erzeugte SysML-v2-Elemente ausweiten. \\
\addlinespace[2pt]
\textbf{SR\_015} & \textbf{Model Content Validation}\par\smallskip The system shall support validation of generated model content before it becomes part of the target architecture. & Absicherung & \textbf{Rationale:} Im Stakeholder Requirement nicht separat dokumentiert.\par\smallskip \textbf{Source:} \texttt{UN\_01\_015; UN\_02\_004; UN\_02\_005} & SF\_009, SF\_010 & \statusopen{Architektur definiert}\par\smallskip Validierungsverantwortung und Zielumgebung sind definiert; generierter Modellinhalt existiert noch nicht. & Strukturelle, semantische und syntaktische Modellvalidierung implementieren. \\
\addlinespace[2pt]
\textbf{SR\_034} & \textbf{Insufficient Evidence Identification}\par\smallskip The system shall identify generated engineering results that are insufficiently supported by available source information. & Absicherung & \textbf{Rationale:} Im Stakeholder Requirement nicht separat dokumentiert.\par\smallskip \textbf{Source:} \texttt{UN\_01\_028} & SF\_002, SF\_004, SF\_007 & \statusdone{Implementiert und verifiziert}\par\smallskip Gaps, Coverage Findings und Attention Conditions kennzeichnen unzureichend gestützte Ergebnisse. & Evidenzschwellen für spätere Modellkandidaten operationalisieren. \\
\addlinespace[2pt]
\textbf{SR\_035} & \textbf{Human Review Decision Association}\par\smallskip The system shall associate each human review decision with the reviewed result. & Absicherung & \textbf{Rationale:} Im Stakeholder Requirement nicht separat dokumentiert.\par\smallskip \textbf{Source:} \texttt{UN\_04\_008} & SF\_007, SF\_011 & \statusdone{Implementiert und verifiziert}\par\smallskip Human Review Decisions sind an Zielinhalt und Validation Fingerprint gebunden und persistent gespeichert. & Association auf modellbezogene Review Decisions erweitern. \\
\addlinespace[2pt]
\textbf{SR\_036} & \textbf{Approved Information Boundary}\par\smallskip The system shall prevent unapproved intermediate results from being treated as approved engineering information. & Absicherung & \textbf{Rationale:} Processing evidence and preliminary results shall remain distinct from approved engineering knowledge.\par\smallskip \textbf{Source:} \texttt{UN\_04\_009} & SF\_004, SF\_007, SF\_008, SF\_010 & \statusdone{Implementiert und verifiziert}\par\smallskip P4/P9 trennen Kandidaten und Processing Evidence strikt von Approved Engineering Knowledge; P9 endet bei awaiting\_review. & Approved Input Promotion als explizit autorisierten Übergang implementieren. \\
\addlinespace[2pt]
\textbf{SR\_039} & \textbf{Processing History Traceability}\par\smallskip The system shall retain a traceable processing history for each registered engineering source. & Absicherung & \textbf{Rationale:} Im Stakeholder Requirement nicht separat dokumentiert.\par\smallskip \textbf{Source:} \texttt{UN\_04\_011} & SF\_001, SF\_005, SF\_007, SF\_011, SF\_012 & \statusdone{Implementiert und verifiziert}\par\smallskip Immutable Events, Decisions, Artifacts und Aggregationen erhalten eine rekonstruierbare Historie je Source und Run. & Historie bis Approved Input, Model Candidates und CATIA-Synchronisation fortführen. \\
\addlinespace[2pt]
\multicolumn{7}{l}{\cellcolor{CategorySemantic}\textbf{Semantik} -- Konsistente Bedeutung technischer Informationen durch Vokabulare und Ontologien.} \\
\addlinespace[1.5pt]
\textbf{SR\_016} & \textbf{Technical Vocabulary Alignment}\par\smallskip The system shall align extracted information with a standardized technical vocabulary. & Semantik & \textbf{Rationale:} Im Stakeholder Requirement nicht separat dokumentiert.\par\smallskip \textbf{Source:} \texttt{UN\_01\_016; UN\_04\_005} & SF\_003 & \statusdone{Implementiert und verifiziert}\par\smallskip Turing Core Vocabulary, Ontology Registry, Glossary Candidates und Terminology Mappings unterstützen die technische Vokabularausrichtung. & Domänenspezifische Erweiterungen und Evaluationsmetriken ergänzen. \\
\addlinespace[2pt]
\textbf{SR\_017} & \textbf{Semantic Equivalence Handling}\par\smallskip The system shall interpret semantically equivalent information consistently across heterogeneous source data. & Semantik & \textbf{Rationale:} Im Stakeholder Requirement nicht separat dokumentiert.\par\smallskip \textbf{Source:} \texttt{UN\_01\_017; UN\_02\_007} & SF\_003, SF\_011 & \statusdone{Implementiert und verifiziert}\par\smallskip Terminology Mapping, Referenzkonzepte und Human Decisions ermöglichen konsistente Behandlung semantischer Äquivalenz. & Konsistenz über spätere Architektur- und Modellartefakte nachweisen. \\
\addlinespace[2pt]
\textbf{SR\_018} & \textbf{Ontology Conflict Detection}\par\smallskip The system shall identify ontology conflicts before the affected information is used for downstream engineering processing. & Semantik & \textbf{Rationale:} Im Stakeholder Requirement nicht separat dokumentiert.\par\smallskip \textbf{Source:} \texttt{UN\_01\_018; UN\_04\_004} & SF\_003 & \statuspartial{Teilweise implementiert}\par\smallskip Ontologien sind versioniert, referenzvalidiert und integrity-checked; Konflikte können als semantische Findings sichtbar werden. & Explizite Ontology-Conflict-Klassifikation und downstream blocking vervollständigen. \\
\addlinespace[2pt]

\end{longtable}
\normalsize
\end{landscape}
\pagestyle{fancy}

\subsection{Zusammenfassung der Requirement-Abdeckung}

Alle 39 Stakeholder Requirements sind im aktuellen Systemmodell mindestens
einer Systemfunktion zugeordnet. Die Implementierungsbewertung ergibt für die
Phase-F/P-Baseline folgenden Stand:

\begin{table}[htbp]
\centering
\caption{Zusammenfassung des Implementierungsstands auf Stakeholder-Requirement-Ebene}
\label{tab:sr_status_summary}
\begin{tabular}{lcc}
\toprule
\textbf{Status} & \textbf{Anzahl} & \textbf{Anteil} \\
\midrule
Implementiert und verifiziert & 16 & 41.0\,\% \\
Teilweise implementiert & 9 & 23.1\,\% \\
Architektur definiert & 13 & 33.3\,\% \\
Optionaler Ausblick & 1 & 2.6\,\% \\
\midrule
\textbf{Gesamt} & \textbf{39} & \textbf{100\,\%} \\
\bottomrule
\end{tabular}
\end{table}

Der aktuelle Implementierungsschwerpunkt liegt auf der kontrollierten
Verarbeitung textueller Bestandsinformationen: Projekt- und Quellenverwaltung,
Ingestion, Ambiguitätserkennung, semantische Zuordnung, Coverage, Human-Review-
Evidenz, Traceability und Statusdarstellung sind weit fortgeschritten oder
bereits verifiziert. Der größte offene Block betrifft die modellbildenden
Requirements: Architekturableitung, Modellintegration, strukturelle und
semantische Validierung sowie die automatische SysML-v2-Erzeugung.

\begin{tcolorbox}[
  enhanced,
  breakable,
  colback=gray!4,
  colframe=TuringSlate,
  title={Hinweis für die finale Fassung der Masterarbeit},
  fonttitle=\bfseries,
  arc=2mm
]
Nach Abschluss des Systems sollte die Spalte \emph{Aktueller Stand} erneut
gegen die finale Implementierung, Testevidenz und Evaluation geprüft werden.
Insbesondere die Requirements SR\_006 bis SR\_010, SR\_015, SR\_020 bis
SR\_025 sowie SR\_027 und SR\_028 werden sich durch die späteren
Architektur- und SysML-v2-Funktionen voraussichtlich wesentlich verändern.
\end{tcolorbox}


\clearpage
\section*{Verwendete Arbeitsgrundlagen}
\addcontentsline{toc}{section}{Verwendete Arbeitsgrundlagen}
\begin{itemize}[leftmargin=7mm,itemsep=2mm]
  \item Autoritatives CATIA-SysML-v2-Modell mit 39 Stakeholder Requirements,
  102 System Requirements und 12 Systemfunktionen.
  \item Transitive Traceability zwischen Stakeholder Requirements,
  System Requirements und Systemfunktionen.
  \item Verifizierte Phase-F/P-Implementierungsbaseline,
  Commit \texttt{26acace4\ldots c7705}, 3.808 bestandene Tests und
  manueller P9-Akzeptanzaudit.
  \item Feature- und Requirement-Coverage-Matrix,
  Arbeitsstand 31.07.2026.
\end{itemize}

\end{document}
